\section*{Capítulo \ref{ch:g12:light-as-wave} --- A luz como onda: difração e interferência}

\begin{solution}{exo:g12:light-as-wave:1}
$f = c/\lambda$: vermelho, $\num{3.00e8}/\num{7.0e-7} \approx
\qty{4.3e14}{Hz}$; violeta, $\approx \qty{7.5e14}{Hz}$ --- o violeta é
maior (\omterm{def:g12:mechanical-waves:wavelength}{comprimento de onda} menor, mesma velocidade $c$).
\end{solution}

\begin{solution}{exo:g12:light-as-wave:2}
$\theta = \num{5.32e-7}/\num{2.0e-4} = \qty{2.7e-3}{rad}$;
$L = 2\theta D = 2 \times \num{2.66e-3} \times 1.5 \approx
\qty{8.0}{mm}$.
\end{solution}

\begin{solution}{exo:g12:light-as-wave:3}
Voz: $\lambda = 340/425 = \qty{0.80}{m}$, logo $\lambda/a = 1.0$ ---
\omterm{def:g12:light-as-wave:diffraction}{difração} máxima, e o \omterm{def:g10:signals-and-waves:sound}{som} inunda o corredor. Luz:
$\lambda/a \approx \num{5.5e-7}/0.80 \approx \num{7e-7}$ ---
espalhamento desprezível: a luz segue reta e a esquina a esconde.
\end{solution}

\begin{solution}{exo:g12:light-as-wave:4}
Em $M$: $\delta = \qty{6.0}{cm} = 3\lambda$, \omterm{prop:g12:light-as-wave:conditions}{construtiva} --- \omterm{def:g10:signals-and-waves:sound}{som} forte.
Em $N$: $\delta = \qty{5.0}{cm} = (2 + \tfrac12)\lambda$, \omterm{prop:g12:light-as-wave:conditions}{destrutiva} ---
silêncio (ou quase).
\end{solution}

\begin{solution}{exo:g12:light-as-wave:5}
$i = \lambda D/a = \num{6.33e-7} \times 1.8/\num{2.5e-4} \approx
\qty{4.6}{mm}$. Dentro de $\pm\qty{1.0}{cm}$: $|k| \leq 10/4.6 = 2.2$,
logo $k = -2, \dots, 2$: cinco franjas claras.
\end{solution}

\begin{solution}{exo:g12:light-as-wave:6}
$d = 2\lambda D/L = 2 \times \num{6.5e-7} \times 1.6/0.028 \approx
\qty{74}{\micro m}$.
\end{solution}

\begin{solution}{exo:g12:light-as-wave:7}
$i = \qty{4.7}{mm}$, logo
$\lambda = i\,a/D = \num{4.7e-3} \times \num{2.0e-4}/1.5 \approx
\qty{6.3e-7}{m} = \qty{630}{nm}$: vermelho. Medir dez \omterm{def:g12:light-as-wave:fringes}{interfranjas}
divide por dez o erro de leitura da régua.
\end{solution}

\begin{solution}{exo:g12:light-as-wave:8}
As luminárias têm \omterm{def:g10:signals-and-waves:frequency}{frequências} (grosso modo) compatíveis, mas \emph{não
estão em sincronia}: cada uma emite trens de \omterm{def:g10:signals-and-waves:wave}{onda} curtos com fase que
salta ao acaso, de modo que as fontes são incoerentes e as franjas
instantâneas se deslocam bilhões de vezes por segundo --- e o olho vê a
média, uniforme. Young ilumina as duas fendas com \emph{uma única} \omterm{def:g10:signals-and-waves:wave}{onda}:
as duas cópias herdam juntas cada salto de fase e permanecem em
sincronia.
\end{solution}

\begin{solution}{exo:g12:light-as-wave:9}
Da primeira à sétima: $6i = \qty{18.0}{mm}$, logo $i = \qty{3.0}{mm}$;
$\lambda = i\,a/D = \num{3.0e-3} \times \num{3.0e-4}/1.4 \approx
\qty{6.4e-7}{m} \approx \qty{640}{nm}$.
\end{solution}

\begin{solution}{exo:g12:light-as-wave:10}
$\sin\theta = k\lambda/d = 0.396\,k$: para $k = 1$,
$\theta = 23^\circ$; para $k = 2$, $\sin\theta = 0.79$,
$\theta = 52^\circ$; $k = 3$ exigiria $\sin\theta = 1.19 > 1$:
impossível. Maior ordem: $k = 2$.
\end{solution}

\begin{solution}{exo:g12:light-as-wave:11}
$\theta = \num{5.5e-7}/0.10 = \qty{5.5e-6}{rad}$. Os faróis subentendem
$1.5/\num{1.0e5} = \qty{1.5e-5}{rad}$, cerca de três vezes o borrão:
separados, por pouco.
\end{solution}

\begin{solution}{exo:g12:light-as-wave:12}
As reflexões das faces da frente e de trás \omterm{prop:g12:light-as-wave:conditions}{interferem}; em \omterm{def:g10:light-spectra:white}{luz branca} cada
espessura reforça alguns \omterm{def:g12:mechanical-waves:wavelength}{comprimentos de onda} e cancela outros, daí uma
cor. A espessura da película depende só da altura, de modo que os pontos
de mesma cor formam faixas horizontais. O escorrimento afina a película,
e por isso a espessura que pintava uma dada cor fica cada vez mais
embaixo: as faixas descem devagar.
\end{solution}

\begin{solution}{exo:g12:light-as-wave:13}
$i = \lambda D/a$: violeta,
$\num{4.0e-7} \times 1.5/\num{2.0e-4} = \qty{3.0}{mm}$; vermelho,
\qty{5.6}{mm}. No centro: todas as cores claras --- uma franja branca. A
poucos milímetros dali os padrões de cores já se descolaram: franjas
iridescentes. Mais longe, os máximos de todas as cores se superpõem em
toda parte: um borrão branco uniforme.
\end{solution}

\begin{solution}{exo:g12:light-as-wave:14}
$\theta = \num{5.5e-7}/2.4 = \qty{2.3e-7}{rad}$; a \qty{250}{km} isso dá
$\num{2.3e-7} \times \num{2.5e5} \approx \qty{6}{cm}$. Ele consegue
contar carros e localizar uma pessoa, mas letra de jornal (com
milímetros) fica vinte vezes abaixo do limite de \omterm{def:g12:light-as-wave:diffraction}{difração}: a afirmação é
lenda.
\end{solution}

\begin{solution}{exo:g12:light-as-wave:15}
A \omterm{def:g10:signals-and-waves:frequency}{frequência} é fixada pela fonte: não muda. A velocidade cai para $c/n$,
de modo que $\lambda' = \lambda/n$ e
$i' = i/n = 4.0/1.33 \approx \qty{3.0}{mm}$: as franjas se apertam.
\end{solution}

\begin{solution}{pb:g12:light-as-wave:1}
\textbf{1.} $\theta = \lambda/a = \num{6.33e-7}/\num{1.00e-4} =
\qty{6.3e-3}{rad}$.

\textbf{2.} Semilargura $D\tan\theta \approx D\theta = \lambda D/a$,
logo $L = 2\lambda D/a = 2 \times \num{6.33e-7} \times
3.00/\num{1.00e-4} = \qty{3.8}{cm}$.

\textbf{3.} $(3.8 - 3.7)/3.8 \approx 3\%$: dentro da precisão de uma
régua --- validado.

\textbf{4.} $L \propto 1/a$: reduzir $a$ à metade dobra $L$, logo
$L = \qty{7.6}{cm}$. Obstáculo (ou fenda) mais estreito, padrão mais
largo.

\textbf{5.} $\tan(\num{6.33e-3}) = \num{6.33008e-3}$: eles concordam a
cerca de $\num{1e-5}$ em termos relativos --- a aproximação é muito
melhor do que qualquer medida feita aqui.

\textbf{6.} Pela \cref{prop:g12:light-as-wave:hair}: um obstáculo de
largura $d$ difrata como uma fenda de largura $d$, de modo que
$L = 2\lambda D/d$ continua valendo.

\textbf{7.} $d = 2\lambda D/L = 2 \times \num{6.33e-7} \times
3.00/0.054 \approx \qty{70}{\micro m}$.

\textbf{8.} $L = \qty{5.2}{cm} \to d = \qty{73}{\micro m}$;
$L = \qty{5.6}{cm} \to d = \qty{68}{\micro m}$: logo,
$d = \qty{70 \pm 3}{\micro m}$.

\textbf{9.} $d = \num{3.80e-6}/0.076 = \qty{50}{\micro m}$: o do amigo é
mais fino --- e, como $L \propto 1/d$, o cabelo mais fino lança o padrão
mais largo.

\textbf{10.} Não: uma luminária é branca (cada \omterm{def:g12:mechanical-waves:wavelength}{comprimento de onda} pinta
um padrão diferente, e todos se borram) e incoerente e larga (não há um
feixe único e limpo para difratar). O \omterm{def:g11:color-light-sources:lamps}{laser} é monocromático e
direcional.

\textbf{11.} Só cópias de uma mesma \omterm{def:g10:signals-and-waves:wave}{onda} são \omterm{def:g12:light-as-wave:coherent}{coerentes}; duas fontes
independentes saem de sincronia e as franjas se apagam
(\cref{rem:g12:light-as-wave:lamps}).

\textbf{12.} No centro, $\delta = 0 = 0 \times \lambda$ para todo
$\lambda$: \omterm{prop:g12:light-as-wave:conditions}{construtiva}, qualquer que seja o \omterm{def:g12:mechanical-waves:wavelength}{comprimento de onda}.

\textbf{13.} $i = 4.26/10 = \qty{4.3}{mm}$; uma extensão de dez é lida
com o mesmo erro de régua de uma única \omterm{def:g12:light-as-wave:fringes}{interfranja}, de modo que o erro
sobre $i$ fica dividido por dez.

\textbf{14.} $\lambda = i\,a'/D = \num{4.26e-3} \times
\num{2.50e-4}/2.00 = \num{5.33e-7} \approx \qty{533}{nm}$.

\textbf{15.} \qty{533}{nm} está dentro de $\qty{532 \pm 10}{nm}$:
coerente; e esse \omterm{def:g12:mechanical-waves:wavelength}{comprimento de onda} é verde.

\textbf{16.} $\tan\theta = 0.43/1.00$, logo $\theta = 23^\circ$. Aqui
$\tan\theta = 0.43$ enquanto $\sin\theta = 0.40$: a $23^\circ$ a
identificação de ângulo pequeno entre $\sin$, $\tan$ e $\theta$ erra por
quase $10\%$ --- use as funções exatas.

\textbf{17.} $d = \lambda/\sin\theta = \num{6.33e-7}/0.396 \approx
\qty{1.6}{\micro m}$.

\textbf{18.} $\tan\theta = 1.65$, $\theta = 59^\circ$,
$\sin\theta = 0.855$: $d = \num{6.33e-7}/0.855 \approx
\qty{0.74}{\micro m}$.

\textbf{19.} $1.6/0.74 \approx 2.2$: as trilhas do DVD são duas vezes
mais densas, e os poços dele são correspondentemente menores (lidos com
um \omterm{def:g11:color-light-sources:lamps}{laser} de \omterm{def:g12:mechanical-waves:wavelength}{comprimento de onda} mais curto) --- várias vezes mais dados
no mesmo disco de \qty{12}{cm}.

\textbf{20.} Cabelo \qty{70 \pm 3}{\micro m}; cabelo do amigo
\qty{50}{\micro m}; \omterm{def:g11:color-light-sources:lamps}{laser} verde \qty{533}{nm}; trilhas do CD
\qty{1.6}{\micro m}; trilhas do DVD \qty{0.74}{\micro m} --- um
\omterm{def:g12:mechanical-waves:wavelength}{comprimento de onda} conhecido transforma a geometria de cada padrão num
comprimento, de modo que uma luz de meio micrômetro é uma régua graduada
em meio micrômetro.
\end{solution}
